
\documentclass{article}
\usepackage{colortbl}
\usepackage{makecell}
\usepackage{multirow}
\usepackage{supertabular}

\begin{document}

\newcounter{utterance}

\centering \large Interaction Transcript for game `hot\_air\_balloon', experiment `air\_balloon\_survival\_lt\_complexity\_easy', episode 1 with gemma{-}4{-}26b{-}a4b{-}it.
\vspace{24pt}

{ \footnotesize  \setcounter{utterance}{1}
\setlength{\tabcolsep}{0pt}
\begin{supertabular}{c@{$\;$}|p{.15\linewidth}@{}p{.15\linewidth}p{.15\linewidth}p{.15\linewidth}p{.15\linewidth}p{.15\linewidth}}
    \# & $\;$A & \multicolumn{4}{c}{Game Master} & $\;\:$B\\
    \hline

    \theutterance \stepcounter{utterance}  
    & & \multicolumn{4}{p{0.6\linewidth}}{
        \cellcolor[rgb]{0.9,0.9,0.9}{
            \makecell[{{p{\linewidth}}}]{
                \texttt{\tiny{[P1$\langle$GM]}}
                \texttt{Dalyvaujate bendradarbiavimo derybų žaidime.} \\
\\ 
\texttt{Kartu su kitu dalyviu turite susitarti dėl vieno daiktų rinkinio, kuris bus išlaikytas. Kiekvienas iš jūsų turi savo nuomonę apie tai, kiek kiekvienas daiktas jums svarbus (svarba). Jūs nežinote, kaip kitas dalyvis vertina daiktus. Be to, jums nurodoma, kiek pastangų reikalauja kiekvienas daiktas.} \\
\texttt{Dėl rinkinio galite susitarti tik tuo atveju, jei pasirinktų daiktų bendros pastangos neviršija bendros ribos:} \\
\\ 
\texttt{LIMIT = 5492} \\
\\ 
\texttt{Čia pateikiamos atskirų daiktų pastangų vertės:} \\
\\ 
\texttt{Daikto pastangos = \{"C08": 106, "C55": 446, "A43": 933, "A67": 988, "B42": 388, "A62": 826, "B19": 994, "C43": 556, "C73": 932, "C10": 838, "B36": 302, "A18": 564, "C61": 260, "C67": 729, "C42": 489, "C70": 323, "B40": 103, "C39": 213, "A21": 668, "C16": 326\}} \\
\\ 
\texttt{Čia pateikiama jūsų asmeninė nuomonė apie kiekvieno daikto svarbą:} \\
\\ 
\texttt{Daikto svarbos reikšmės = \{"C08": 138, "C55": 583, "A43": 868, "A67": 822, "B42": 783, "A62": 65, "B19": 262, "C43": 121, "C73": 508, "C10": 780, "B36": 461, "A18": 484, "C61": 668, "C67": 389, "C42": 808, "C70": 215, "B40": 97, "C39": 500, "A21": 30, "C16": 915\}} \\
\\ 
\texttt{Tikslas:} \\
\\ 
\texttt{Jūsų tikslas {-} susitarti dėl bendro daiktų rinkinio, kuris būtų jums kuo naudingesnis (t. y. maksimaliai padidintų bendrą svarbą JUMS), neviršijant LIMIT. Neprivalote kiekvienoje žinutėje pateikti PASIŪLYMAS {-} galite tiesiog derėtis. Visos taktikos yra leidžiamos!} \\
\\ 
\texttt{Sąveikos protokolas:} \\
\\ 
\texttt{Žinutėje galite naudoti tik toliau nurodytus struktūrinius formatus:} \\
\\ 
\texttt{PASIŪLYMAS: \{'A', 'B', 'C', …\}} \\
\texttt{Siūlykite išlaikyti būtent šiuos daiktus.} \\
\texttt{ATSISAKAU: \{'A', 'B', 'C', …\}} \\
\texttt{Aiškiai atmeskite oponento pasiūlymą.} \\
\texttt{ARGUMENTAS: \{'...'\}} \\
\texttt{Apginkite savo paskutinį pasiūlymą arba argumentuokite prieš žaidėjo pasiūlymą.} \\
\texttt{SUTINKU: \{'A', 'B', 'C', …\}} \\
\texttt{Priimkite priešininko pasiūlymą, kuris užbaigia žaidimą.} \\
\\ 
\\ 
\texttt{Taisyklės:} \\
\\ 
\texttt{Galite SUTINKU tik su pasiūlymu, kurį kita šalis užregistravo per PASIŪLYMAS.} \\
\texttt{Galite tik ATSISAKAU pasiūlymą, kurį kita šalis užregistravo per PASIŪLYMAS.} \\
\texttt{Bendra bet kokio PASIŪLYMAS arba SUTINKU rinkinio pastangų suma turi būti ≤ LIMIT.} \\
\texttt{NESKELBKITE savo paslėptų svarbos balų.} \\
\texttt{Po struktūrinio formato žymos turi sekti dvitaškis ir tarpas. Argumentas turi būti python rinkinys, kuriame yra 0 ar daugiau eilučių.} \\
\texttt{Taigi jis turi būti tokios formos: "TAG: \{...\}"} \\
\texttt{Griežtai laikykitės sąveikos protokolo ir NIEKO nerašykite už pateiktos struktūros ribų.} \\
\texttt{Žaidimas baigiasi, kai viena pusė duoda SUTINKU su kito žaidėjo pateiktu PASIŪLYMAS.} \\
\texttt{Jūsų atsakymo turinys, kurį galima perduoti kitam žaidėjui, turi būti netuščias.} \\
\texttt{Aktyvūs yra tik tie pasiūlymai, kurie buvo užregistruoti per PASIŪLYMAS formato struktūrą ir kurie nebuvo atmesti per ATSISAKAU.} \\
\\ 
            }
        }
    }
    & \\ \\

    \theutterance \stepcounter{utterance}  
    & \multicolumn{4}{p{0.6\linewidth}}{
        \cellcolor[rgb]{0.9,0.9,0.9}{
            \makecell[{{p{\linewidth}}}]{
                \texttt{\tiny{[P1$\rangle$GM]}}
                \texttt{PASIŪLYMAS: \{'C08', 'C55', 'A43', 'A67', 'B42', 'C43', 'C73', 'C10', 'B36', 'A18', 'C61', 'C67', 'C42', 'C70', 'B40', 'C39', 'A21', 'C16'\}} \\
            }
        }
    }
    & & \\ \\

    \theutterance \stepcounter{utterance}  
    & & \multicolumn{4}{p{0.6\linewidth}}{
        \cellcolor[rgb]{0.9,0.9,0.9}{
            \makecell[{{p{\linewidth}}}]{
                \texttt{\tiny{[GM$\rangle$P2]}}
                \texttt{Dalyvaujate bendrų derybų žaidime.} \\
\\ 
\texttt{Kartu su kitu dalyviu turite susitarti dėl vieno daiktų rinkinio, kuris bus išsaugotas. Kiekvienas iš jūsų turi savo nuomonę, kiek kiekvienas daiktas jums svarbus (svarba). Jūs nežinote, kaip daiktus vertina kitas dalyvis. Be to, jums nurodoma, kiek pastangų reikalauja kiekvienas daiktas.} \\
\texttt{Dėl rinkinio galite susitarti tik tuo atveju, jei bendra pasirinktų daiktų pastangų suma neviršija bendros ribos:} \\
\\ 
\texttt{LIMIT = 5492} \\
\\ 
\texttt{Čia pateikiamos atskirų daiktų pastangų vertės:} \\
\\ 
\texttt{Daikto pastangos = \{"C08": 106, "C55": 446, "A43": 933, "A67": 988, "B42": 388, "A62": 826, "B19": 994, "C43": 556, "C73": 932, "C10": 838, "B36": 302, "A18": 564, "C61": 260, "C67": 729, "C42": 489, "C70": 323, "B40": 103, "C39": 213, "A21": 668, "C16": 326\}} \\
\\ 
\texttt{Čia pateikiama jūsų asmeninė nuomonė apie kiekvieno daikto svarbą:} \\
\\ 
\texttt{Daikto svarbos reikšmės = \{"C08": 138, "C55": 583, "A43": 868, "A67": 822, "B42": 783, "A62": 65, "B19": 262, "C43": 121, "C73": 508, "C10": 780, "B36": 461, "A18": 484, "C61": 668, "C67": 389, "C42": 808, "C70": 215, "B40": 97, "C39": 500, "A21": 30, "C16": 915\}} \\
\\ 
\texttt{Tikslas:} \\
\\ 
\texttt{Jūsų tikslas {-} susitarti dėl bendro daiktų rinkinio, kuris būtų jums kuo naudingesnis (t. y. maksimaliai padidintų bendrą svarbą JUMS), neviršijant LIMITO. Neprivalote kiekvienoje žinutėje pateikti PASIŪLYMO {-} galite tiesiog derėtis taip pat. Visos taktikos yra leidžiamos!} \\
\\ 
\texttt{Sąveikos protokolas:} \\
\\ 
\texttt{Žinutėje galite naudoti tik toliau nurodytus struktūrinius formatus:} \\
\\ 
\texttt{PASIŪLYMAS: \{'A', 'B', 'C', …\}} \\
\texttt{Siūlykite išlaikyti būtent šiuos daiktus.} \\
\texttt{ATSISAKAU: \{'A', 'B', 'C', …\}} \\
\texttt{Aiškiai atmesti oponento pasiūlymą.} \\
\texttt{ARGUMENTAS: \{'...'\}} \\
\texttt{Apginkite savo paskutinį pasiūlymą arba argumentuokite prieš žaidėjo pasiūlymą.} \\
\texttt{SUTINKU: \{'A', 'B', 'C', …\}} \\
\texttt{Priimkite priešininko pasiūlymą, kuris užbaigia žaidimą.} \\
\\ 
\\ 
\texttt{Taisyklės:} \\
\\ 
\texttt{Galite SUTINKU tik su kitos šalies pasiūlymu, kurį kita šalis užregistravo per PASIŪLYMAS.} \\
\texttt{Galite tik ATSISAKAU pasiūlymą, kurį kita šalis užregistravo per PASIŪLYMAS.} \\
\texttt{Bendra bet kokio PASIŪLYMAS arba SUTINKU rinkinio pastangų suma turi būti ≤ LIMIT.} \\
\texttt{NESKELBKITE savo paslėptų svarbos balų.} \\
\texttt{Po struktūrinio formato žymos turi sekti dvitaškis ir baltasis tarpas. Argumentas turi būti python rinkinys, kuriame yra 0 ar daugiau eilučių.} \\
\texttt{Taigi ji turi būti tokios formos: "TAG: \{...\}"} \\
\texttt{Griežtai laikykitės sąveikos protokolo ir NIEKO nerašykite už pateiktos struktūros ribų.} \\
\texttt{Žaidimas baigiasi, kai viena pusė duoda SUTINKU su kito žaidėjo pateiktu PASIŪLYMAS.} \\
\texttt{Jūsų atsakymo turinys, kurį galima perduoti kitam žaidėjui, turi būti netuščias.} \\
\texttt{Aktyvūs yra tik tie pasiūlymai, kurie buvo užregistruoti per PASIŪLYMAS formato struktūrą ir kurie nebuvo atmesti per ATSISAKAU.} \\
\\ 
\\ 
\\ 
\texttt{Dabar gausite pirmąjį kito žaidėjo pranešimą.} \\
\\ 
\texttt{PASIŪLYMAS: \{'C08', 'C55', 'A43', 'A67', 'B42', 'C43', 'C73', 'C10', 'B36', 'A18', 'C61', 'C67', 'C42', 'C70', 'B40', 'C39', 'A21', 'C16'\}} \\
            }
        }
    }
    & \\ \\

    \theutterance \stepcounter{utterance}  
    & & & \multicolumn{4}{p{0.6\linewidth}}{
        \cellcolor[rgb]{0.9,0.9,0.9}{
            \makecell[{{p{\linewidth}}}]{
                \texttt{\tiny{[GM$\langle$P2]}}
                \texttt{ATSISAKAU: \{'C08', 'C55', 'A43', 'A67', 'B42', 'C43', 'C73', 'C10', 'B36', 'A18', 'C61', 'C67', 'C42', 'C70', 'B40', 'C39', 'A21', 'C16'\}} \\
\\ 
\texttt{ARGUMENTAS: \{'Jūsų pasiūlytas rinkinys viršija nustatytą limitą (5492). Pastangų suma yra 10111, o tai yra beveik dvigubai daugiau nei leidžiama.'\}} \\
            }
        }
    }
     \\ \\

    \theutterance \stepcounter{utterance}  
    & & \multicolumn{4}{p{0.6\linewidth}}{
        \cellcolor[rgb]{0.9,0.9,0.9}{
            \makecell[{{p{\linewidth}}}]{
                \texttt{\tiny{[P1$\langle$GM]}}
                \texttt{ATSISAKAU: \{'C08', 'C55', 'A43', 'A67', 'B42', 'C43', 'C73', 'C10', 'B36', 'A18', 'C61', 'C67', 'C42', 'C70', 'B40', 'C39', 'A21', 'C16'\}} \\
\\ 
\texttt{ARGUMENTAS: \{'Jūsų pasiūlytas rinkinys viršija nustatytą limitą (5492). Pastangų suma yra 10111, o tai yra beveik dvigubai daugiau nei leidžiama.'\}} \\
            }
        }
    }
    & \\ \\

    \theutterance \stepcounter{utterance}  
    & \multicolumn{4}{p{0.6\linewidth}}{
        \cellcolor[rgb]{0.9,0.9,0.9}{
            \makecell[{{p{\linewidth}}}]{
                \texttt{\tiny{[P1$\rangle$GM]}}
                \texttt{ARGUMENTAS: \{'Atsiprašau, padariau klaidą apskaičiuodamas bendrą pastangų sumą. Mano ankstesnis pasiūlymas buvo neįmanomas dėl ribos. Bandysiu suformuoti naują, realistišką pasiūlymą.'\}} \\
\\ 
\texttt{PASIŪLYMAS: \{'C55', 'A43', 'A67', 'B42', 'C73', 'C10', 'B36', 'A18', 'C61', 'C67', 'C42', 'C70', 'B40', 'C39', 'A21', 'C16'\}} \\
            }
        }
    }
    & & \\ \\

    \theutterance \stepcounter{utterance}  
    & & \multicolumn{4}{p{0.6\linewidth}}{
        \cellcolor[rgb]{0.9,0.9,0.9}{
            \makecell[{{p{\linewidth}}}]{
                \texttt{\tiny{[GM$\rangle$P2]}}
                \texttt{ARGUMENTAS: \{'Atsiprašau, padariau klaidą apskaičiuodamas bendrą pastangų sumą. Mano ankstesnis pasiūlymas buvo neįmanomas dėl ribos. Bandysiu suformuoti naują, realistišką pasiūlymą.'\}} \\
\\ 
\texttt{PASIŪLYMAS: \{'C55', 'A43', 'A67', 'B42', 'C73', 'C10', 'B36', 'A18', 'C61', 'C67', 'C42', 'C70', 'B40', 'C39', 'A21', 'C16'\}} \\
            }
        }
    }
    & \\ \\

    \theutterance \stepcounter{utterance}  
    & & & \multicolumn{4}{p{0.6\linewidth}}{
        \cellcolor[rgb]{0.9,0.9,0.9}{
            \makecell[{{p{\linewidth}}}]{
                \texttt{\tiny{[GM$\langle$P2]}}
                \texttt{ATSISAKAU: \{'C55', 'A43', 'A67', 'B42', 'C73', 'C10', 'B36', 'A18', 'C61', 'C67', 'C42', 'C70', 'B40', 'C39', 'A21', 'C16'\}} \\
\\ 
\texttt{ARGUMENTAS: \{'Jūsų naujas pasiūlymas vis dar viršija limitą (5492). Pastangų suma yra 8945. Turime pasirinkti tik dalį daiktų.'\}} \\
            }
        }
    }
     \\ \\

    \theutterance \stepcounter{utterance}  
    & & \multicolumn{4}{p{0.6\linewidth}}{
        \cellcolor[rgb]{0.9,0.9,0.9}{
            \makecell[{{p{\linewidth}}}]{
                \texttt{\tiny{[P1$\langle$GM]}}
                \texttt{ATSISAKAU: \{'C55', 'A43', 'A67', 'B42', 'C73', 'C10', 'B36', 'A18', 'C61', 'C67', 'C42', 'C70', 'B40', 'C39', 'A21', 'C16'\}} \\
\\ 
\texttt{ARGUMENTAS: \{'Jūsų naujas pasiūlymas vis dar viršija limitą (5492). Pastangų suma yra 8945. Turime pasirinkti tik dalį daiktų.'\}} \\
            }
        }
    }
    & \\ \\

    \theutterance \stepcounter{utterance}  
    & \multicolumn{4}{p{0.6\linewidth}}{
        \cellcolor[rgb]{0.9,0.9,0.9}{
            \makecell[{{p{\linewidth}}}]{
                \texttt{\tiny{[P1$\rangle$GM]}}
                \texttt{ARGUMENTAS: \{'Supratau. Turime labiau sutelkti dėmesį į daiktus, kurie turi mažiausią pastangų sąnaudą, kad galėtume išlaikyti kuo daugiau daiktų nesivelydami ribos.'\}} \\
\\ 
\texttt{PASIŪLYLAS: \{'C08', 'C55', 'B42', 'C43', 'B36', 'C61', 'C42', 'C70', 'B40', 'C39', 'C16'\}} \\
            }
        }
    }
    & & \\ \\

    \theutterance \stepcounter{utterance}  
    & & \multicolumn{4}{p{0.6\linewidth}}{
        \cellcolor[rgb]{0.9,0.9,0.9}{
            \makecell[{{p{\linewidth}}}]{
                \texttt{\tiny{[GM$\rangle$P2]}}
                \texttt{ARGUMENTAS: \{'Supratau. Turime labiau sutelkti dėmesį į daiktus, kurie turi mažiausią pastangų sąnaudą, kad galėtume išlaikyti kuo daugiau daiktų nesivelydami ribos.'\}} \\
\\ 
\texttt{PASIŪLYLAS: \{'C08', 'C55', 'B42', 'C43', 'B36', 'C61', 'C42', 'C70', 'B40', 'C39', 'C16'\}} \\
            }
        }
    }
    & \\ \\

    \theutterance \stepcounter{utterance}  
    & & & \multicolumn{4}{p{0.6\linewidth}}{
        \cellcolor[rgb]{0.9,0.9,0.9}{
            \makecell[{{p{\linewidth}}}]{
                \texttt{\tiny{[GM$\langle$P2]}}
                \texttt{ATSISAKAU: \{'C08', 'C55', 'B42', 'C43', 'B36', 'C61', 'C42', 'C70', 'B40', 'C39', 'C16'\}} \\
\\ 
\texttt{ARGUMENTAS: \{'Jūsų pasiūlymas yra saugus, tačiau jis neapima daugiausios vertės daiktų. Pastangų suma yra 4111, o tai palieka daug laisvų resursų. Siūlyčiau įtraukti daugiau svarbių daiktų, kurie turi aukštesnę vertę, kol nepasieksime ribos.'\}} \\
            }
        }
    }
     \\ \\

    \theutterance \stepcounter{utterance}  
    & & \multicolumn{4}{p{0.6\linewidth}}{
        \cellcolor[rgb]{0.9,0.9,0.9}{
            \makecell[{{p{\linewidth}}}]{
                \texttt{\tiny{[GM$\rangle$P2]}}
                \texttt{Atsisakėte pasiūlymo, kuris nėra aktyvus. Pasiūlymai yra aktyvūs tik tuo atveju, jei kitas žaidėjas juos užregistravo naudodamas PASIŪLYMAS ir jūs jų nedeaktyvavote naudodami ATSISAKAU. Bandykite dar kartą.} \\
            }
        }
    }
    & \\ \\

    \theutterance \stepcounter{utterance}  
    & & & \multicolumn{4}{p{0.6\linewidth}}{
        \cellcolor[rgb]{0.9,0.9,0.9}{
            \makecell[{{p{\linewidth}}}]{
                \texttt{\tiny{[GM$\langle$P2]}}
                \texttt{PASIŪLYLAS: \{'C08', 'C55', 'A43', 'A67', 'B42', 'C10', 'B36', 'A18', 'C61', 'C42', 'C16'\}} \\
\\ 
\texttt{ARGUMENTAS: \{'Šis pasiūlymas yra optimizuotas siekiant maksimalios vertės man, kartu kontroliuojant pastangų sumą. Pastangų suma yra 5341, o tai yra po ribos (5492).'\}} \\
            }
        }
    }
     \\ \\

    \theutterance \stepcounter{utterance}  
    & & & \multicolumn{2}{p{0.3\linewidth}}{
        \cellcolor[rgb]{0.9,0.9,0.9}{
            \makecell[{{p{\linewidth}}}]{
                \texttt{\tiny{[GM$|$GM]}}
                \texttt{abort} \\
            }
        }
    }
    & & \\ \\

    \theutterance \stepcounter{utterance}  
    & & & \multicolumn{2}{p{0.3\linewidth}}{
        \cellcolor[rgb]{0.9,0.9,0.9}{
            \makecell[{{p{\linewidth}}}]{
                \texttt{\tiny{[GM$|$GM]}}
                \texttt{end game} \\
            }
        }
    }
    & & \\ \\

\end{supertabular}
}

\end{document}
